\documentclass{article}
\usepackage{amsmath}
\usepackage{amssymb}
\usepackage{amsfonts}
\usepackage{graphicx}
\usepackage{fn2end} % To place footnotes at the end of the chapter/section if possible, but standard placement here.
\usepackage{booktabs} % Not strictly needed, but good practice if tables were present

\begin{document}
\pagestyle{myheadings}
\markboth{Prefaces}{Preface to the Second English Edition}

\begin{center}
{\Large \bfseries Prefaces}
\end{center}

\vspace{1cm}

{\large \bfseries Preface to the Second English Edition}

Science has not stood still in the years since the first English edition of this book
was published. For example, Fermat's last theorem has been proved, the Poincaré
conjecture is now a theorem, and the Higgs boson has been discovered. Other events
in science, while not directly related to the contents of a textbook in classical math-
ematical analysis, have indirectly led the author to learn something new, to think
over something familiar, or to extend his knowledge and understanding. All of this
additional knowledge and understanding end up being useful even when one speaks
about something apparently completely unrelated.\footnote{There is a story about Erd\H{o}s, who, like Hadamard, lived a very long mathematical and human life. When he was quite old, a journalist who was interviewing him asked him about his age. Erd\H{o}s replied, after deliberating a bit, ``I remember that when I was very young, scientists established that the Earth was two billion years old. Now scientists assert that the Earth is four and a half billion years old. So, I am approximately two and a half billion years old.''}

In addition to the original Russian edition, the book has been published in En-
glish, German, and Chinese. Various attentive multilingual readers have detected
many errors in the text. Luckily, these are local errors, mostly misprints. They have
assuredly all been corrected in this new edition.

But the main difference between the second and first English editions is the addi-
tion of a series of appendices to each volume. There are six of them in the first and
five of them in the second. So as not to disturb the original text, they are placed at the
end of each volume. The subjects of the appendices are diverse. They are meant to be
useful to students (in mathematics and physics) as well as to teachers, who may be
motivated by different goals. Some of the appendices are surveys, both prospective
and retrospective. The final survey contains the most important conceptual achieve-
ments of the whole course, which establish connections between analysis and other
parts of mathematics as a whole.

\vfill
\begin{flushright}
V
\end{flushright}
\end{document}